% =============================================================================
%  2.3 Рендеринг
% =============================================================================
\section{Рендеринг}

Процесс рендеринга в React состоит из двух этапов:
\term{согласование}{Reconciliation} (вычисление разницы) и
\term{фиксация}{Commit} (применение изменений к DOM).
Понимание этого разделения критически важно для оптимизации
производительности приложений.

React использует технику \term{пакетирования обновлений}{Batching},
группируя несколько вызовов \texttt{setState} в один цикл рендеринга.

\subsection{Virtual DOM и диффинг}

\subsubsection{Эвристики алгоритма}

Алгоритм сравнения основан на двух эвристиках: элементы разных типов
порождают разные деревья, а атрибут \texttt{key} указывает, какие
дочерние элементы стабильны между рендерами.

Без ключей сложность переупорядочивания списка из $n$ элементов
может достигать $O(n^2)$; с ключами — $O(n)$:
\begin{equation}
  C_{\text{reorder}} = \begin{cases}
    O(n)   & \text{с ключами,}\\
    O(n^2) & \text{без ключей.}
  \end{cases}
\end{equation}

\begin{code}{jsx}
function TodoList({ items }) {
  return (
    <ul>
      {items.map(item => (
        <li key={item.id}>{item.text}</li>
      ))}
    </ul>
  );
}
\end{code}

\note{Атрибут \texttt{key} должен быть стабильным, уникальным и
предсказуемым.  Индекс массива подходит только для статических
списков, которые не переупорядочиваются и не фильтруются.}

\important{Не используйте \texttt{Math.random()} или текущее время
в качестве \texttt{key} — это приводит к полному пересозданию элементов
при каждом рендере и потере состояния дочерних компонентов.}
